\documentclass[a4paper,11pt]{article}
\usepackage{exam-style-842-unit}
\examinehead{842 信息化概论 · 第07单元练习册}{练习二 参考答案与解析}
\begin{document}

\unitanswerscover{842 信息化概论 · 单元练习册}{第07单元参考答案}{练习二（人工智能应用与发展）的答案与解析；每道判断题均给出改正后的正确表述，教材未展开处已注明来源。}

\answerpart{一、单项选择题（10 题，每题 3 分，共 30 分）}

\q C。\begin{analysis}教材 9.6 开篇指出，人工智能技术可分为符号智能和计算智能，其应用十分广泛，主要涉及 15 个方面的应用（教材 9.6，印刷页 155）。A、B、D 均不是教材给出的数目。\end{analysis}

\q A。\begin{analysis}自动规划、调度与配置问题都属于人工智能经典问题之一的约束满足问题，这类问题的解决体现了计算机的创造性（教材 9.6.2，印刷页 156）。B 分类、C 聚类、D 降维都是机器学习的常见任务类型，不是教材对规划调度配置问题的归类。\end{analysis}

\q B。\begin{analysis}机器定理证明的方法主要有自然演绎法、判定法、定理证明器、计算机辅助证明四类。其中判定法是对一类问题找出统一的在计算机上可实现的算法解，著名成果是我国数学家吴文俊教授 1977 年提出的初等几何定理证明方法（教材 9.6.3，印刷页 156）。A 自然演绎法依据推理规则从前提和公理推出定理；C 定理证明器研究一切可判定问题的证明方法；D 计算机辅助证明以计算机为辅助工具完成大量计算、推理和穷举。\end{analysis}

\q B。\begin{analysis}智能控制具有两个显著特点：一是同时具有知识表示的非数学广义世界模型和传统数学模型混合表示的控制过程；二是其核心在于高层控制，即组织级控制，任务是对实际环境或过程进行组织，即决策与规划，以实现广义问题求解（教材 9.6.6，印刷页 157）。A、C、D 都是智能控制实施中的具体条件，不是教材所说的“核心”。\end{analysis}

\q C。\begin{analysis}智能决策支持系统是 20 世纪 80 年代初由决策支持系统与人工智能特别是专家系统相结合产生的，它既发挥传统决策支持系统数值分析的优势，又发挥专家系统知识及知识处理的特长，既可定量分析又可定性分析，能有效处理半结构化和非结构化问题（教材 9.6.8，印刷页 158）。A、B 都只强调一个方面，D 否认了与专家系统的关系，均错。\end{analysis}

\q B。\begin{analysis}教材 9.6.12（印刷页 159）指出，模式识别是指利用计算机进行物体识别，这里的物体一般指文字、符号、图形、图像、语音、声音及传感器信息等形式的实体对象，并不包括概念、思想、意识等抽象或虚拟对象。A 恰好说反，C、D 把物体范围限得过窄。\end{analysis}

\q A。\begin{analysis}教材 9.6.13（印刷页 159）指出：数据挖掘（DM，也称数据开采、数据采掘等）和数据库中的知识发现（KDD）的本质含义是一样的，只是前者主要用于统计、数据分析、数据库和信息系统等领域，后者则主要用于人工智能和机器学习等领域。B 否认联系，C、D 把二者各自的应用领域说反或说窄。\end{analysis}

\q B。\begin{analysis}1997 年，IBM 的“深蓝”计算机以 2 胜 3 平 1 负的战绩击败了蝉联 12 年之久的世界国际象棋冠军加里·卡斯帕罗夫（教材 9.6.14，印刷页 159）。A 沃森是 2011 年参加美国智力问答节目的人工智能程序；C AlphaGo 是 2016 年下围棋的程序；D ENIAC 是世界上早期电子计算机，不是博弈程序。\end{analysis}

\q B。\begin{analysis}2016 年 3 月，阿尔法围棋（AlphaGo）对战世界围棋冠军、职业九段选手李世石，并以 4:1 的总比分获胜；其主要工作原理是“深度学习”，即多层的人工神经网络和训练的方法（教材 9.6.14，印刷页 159—160）。A 专家系统、C 遗传算法、D 蚁群算法都不是 AlphaGo 的主要原理。\end{analysis}

\q A。\begin{analysis}教材 9.6.15（印刷页 160—161）指出，智能机器人一般至少要具备三个要素：一是感觉要素，用来认识周围环境状态；二是运动要素，对外界做出反应性动作；三是思考要素，根据感觉要素得到的信息思考采用什么动作。注意教材 9.2.2 讲智能机器时表述为感觉、反应、思考三要素，两处表述略有差异，本题按 9.6.15 作答。B、C、D 均非教材表述。\end{analysis}

\answerpart{二、多项选择题（3 题，每题 4 分，共 12 分）}

\q ABC。\begin{analysis}机器翻译、模式识别、机器博弈都属于教材 9.6 所列的人工智能应用领域（分别为 9.6.5、9.6.12、9.6.14）。D 错误：人工合成淀粉属于生物技术与合成生物学成果，不是人工智能的应用领域。\end{analysis}

\q ABD。\begin{analysis}A 正确：“深蓝”的超人表现几乎纯粹是靠运算能力，它被输入数百万个国际象棋案例后筛选最佳着法；B 正确：AlphaGo 的主要工作原理是深度学习；D 正确：两者都是机器博弈的典型代表（教材 9.6.14，印刷页 159—160）。C 错误：围棋每回合约有 250 种可能、一盘棋可长达 150 回合，落子可能性远多于国际象棋，即便是运算速度最快的计算机也无法穷举其中一小部分，AlphaGo 靠的是深度学习与两个神经网络的合作，而不是穷举。\end{analysis}

\q ABC。\begin{analysis}A 正确：教材 9.6.15（印刷页 160—161）指出智能机器人一般至少具备感觉、运动、思考三个要素；B 正确：智能机器人的研制几乎需要所有的人工智能技术，是人工智能技术的综合应用；C 正确：教材 9.6.15（印刷页 161）指出工业机器人虽已广泛应用于各大门类工业领域，但主要在结构化环境中执行各类确定性任务。D 错误：教材 9.7.2（印刷页 162—163）明确指出机器人在意识、推理等方面不可能超越人类，不能复制创造力、想象力和直觉。\end{analysis}

\answerpart{三、判断题（6 题，每题 2 分，共 12 分）}

\q √。\begin{analysis}教材 9.6.2（印刷页 156）指出，从问题求解角度看，自动规划、调度、配置三者有一定的内在联系，有时甚至可以互相转化，事实上它们都属于人工智能经典问题之一的约束满足问题。\end{analysis}

\q ×。改正：吴文俊 1977 年提出的初等几何定理证明方法属于判定法，即对一类问题找出统一的在计算机上可实现的算法解。\begin{analysis}教材 9.6.3（印刷页 156）把机器定理证明方法分为自然演绎法、判定法、定理证明器、计算机辅助证明四类，吴文俊方法是判定法的著名成果，不是自然演绎法。\end{analysis}

\q ×。改正：智能决策支持系统既能进行定量分析，又能进行定性分析，可有效处理半结构化和非结构化问题。\begin{analysis}教材 9.6.8（印刷页 158）指出，智能决策支持系统既发挥传统决策支持系统数值分析的优势，又发挥专家系统知识处理的特长，既可定量分析又可定性分析，从而扩大决策支持系统的范围、提高其能力。\end{analysis}

\q ×。改正：模式识别中的“物体”一般指文字、符号、图形、图像、语音、声音及传感器信息等实体对象，不包括概念、思想、意识等抽象或虚拟对象。\begin{analysis}教材 9.6.12（印刷页 159）对模式识别的研究对象作了明确限定，本题把“不包括”说成了“包括”。\end{analysis}

\q √。\begin{analysis}教材 9.6.13（印刷页 159）指出，数据挖掘（DM）和数据库中的知识发现（KDD）的本质含义是一样的，只是前者主要用于统计、数据分析、数据库和信息系统等领域，后者主要用于人工智能和机器学习等领域。\end{analysis}

\q √。\begin{analysis}教材 9.6.15（印刷页 161）指出，智能机器人的研制几乎需要所有的人工智能技术，还涉及其他许多科学技术部门和领域，所以智能机器人是人工智能技术的综合应用，其能力和水平已经成为人工智能技术水平甚至人类科学技术综合水平的一个体现。\end{analysis}

\answerpart{四、填空题（6 空，每空 1 分，共 6 分）}

\q 15。\begin{analysis}教材 9.6 开篇明确“主要涉及以下 15 个方面的应用”（印刷页 155）。\end{analysis}

\q “深蓝”。\begin{analysis}1997 年 IBM 的“深蓝”以 2 胜 3 平 1 负击败国际象棋世界冠军卡斯帕罗夫（教材 9.6.14，印刷页 159）。\end{analysis}

\q 深度学习。\begin{analysis}AlphaGo 的主要工作原理是深度学习，即多层人工神经网络和训练的方法（教材 9.6.14，印刷页 159—160）。\end{analysis}

\q 人工智能。\begin{analysis}智能决策支持系统是决策支持系统与人工智能特别是专家系统相结合的产物（教材 9.6.8，印刷页 158）。\end{analysis}

\q KDD。\begin{analysis}数据挖掘的英文缩写是 DM（Data Mining），数据库中的知识发现的英文缩写是 KDD（Knowledge Discovery in Databases）（教材 9.6.13，印刷页 159）。\end{analysis}

\q 机器情感。\begin{analysis}教材 9.7.3（印刷页 163）指出，未来人工智能可能会向模糊处理、并行化、神经网络和机器情感 4 个方面发展，并认为情感是智能的一部分，赋予计算机情感能力可能是下一个突破。\end{analysis}

\answerpart{五、简答题（4 题，每题 6 分，共 24 分）}

\q 参考答案要点（写出 10 个即可得满分，全部列出不扣分；每个 0.6 分）：
（1）难题求解：求解没有算法解或在现有机器上无法实施、无法完成的困难问题，如旅行商问题、路径规划、调度、预测与诊断。（2）自动规划、调度与配置：设计行动序列、安排任务分配、设计部件组合结构，均属约束满足问题。（3）机器定理证明：用计算机证明数学定理，方法有自然演绎法、判定法、定理证明器和计算机辅助证明。（4）自动程序设计：给出高级描述后由计算机自动生成并验证程序。（5）机器翻译：完全用计算机实现两种语言之间的翻译。（6）智能控制：把人工智能引入控制领域，核心是高层控制即组织级控制。（7）智能管理：把人工智能引入管理领域，建立智能管理系统。（8）智能决策：把人工智能引入决策过程，建立智能决策支持系统。（9）智能通信：在通信网的构建、网管与网控、转接、信息传输与转换等环节实现智能化。（10）智能仿真：用人工智能指导建模、实验运行及结果分析。（11）智能人机接口：将人工智能用于人机交互界面，使人机界面更加灵性化、拟人化、个性化。（12）模式识别：利用计算机识别文字、符号、图形、图像、语音等实体对象。（13）数据挖掘与知识发现：从数据中提取隐含的、有价值的知识。（14）机器博弈：用计算机下棋，经历塞缪尔跳棋、“深蓝”、AlphaGo 等阶段。（15）智能机器人：具备感觉、运动、思考三要素，是人工智能技术的综合应用。（教材 9.6，印刷页 155—161）

\q 参考答案要点：
（1）含义（2 分）：智能决策支持系统是决策支持系统与人工智能特别是专家系统相结合的产物，它既发挥传统决策支持系统数值分析的优势，又发挥专家系统知识及知识处理的特长，既可定量分析又可定性分析，能有效解决半结构化和非结构化问题（教材 9.6.8，印刷页 158）。
（2）智能部件的模式（2 分）：智能部件可以有多种模式，如专家系统模式、知识库系统模式等。专家系统模式把专家系统作为智能部件，是目前比较流行的一种模式，适合以知识处理为主的问题，但与决策支持系统的接口比较困难；知识库系统模式以知识库作为智能部件，此时决策支持系统由模型库、方法库、数据库、知识库组成四库系统，接口较易实现，整体性能也较好。
（3）与传统决策支持系统的区别（2 分）：传统决策支持系统以数值分析为主、擅长定量分析、主要处理结构化问题；智能决策支持系统在传统决策支持系统基础上加上智能部件，实现定量与定性分析结合，扩大了决策支持系统的范围、提高了决策支持系统的能力，可处理半结构化和非结构化问题。

\q 参考答案要点：
（1）联系（2 分）：数据挖掘（DM，也称数据开采、数据采掘等）和数据库中的知识发现（KDD）的本质含义是一样的，都是指从大量数据中提取隐含的、有价值的知识，二者在文献中一般同时列出，并已成为人工智能应用的一个热门领域和研究方向（教材 9.6.13，印刷页 159）。
（2）区别（2 分）：DM 主要用于统计、数据分析、数据库和信息系统等领域；KDD 则主要用于人工智能和机器学习等领域，二者的侧重点和学科归属不同。其涉及范围非常广泛，如企业数据、商业数据、科学实验数据、管理决策数据、Web 数据的挖掘和发现。
（3）5 条主线（2 分）：教材 9.6 末（印刷页 161）归纳了人工智能技术形成的 5 条主线：一是从专家（知识）系统到代理人（Agent）系统和智能机器人系统；二是从机器学习到数据挖掘和数据库中的知识发现；三是从基于图搜索的问题求解到基于各种智能算法的问题求解；四是从单机环境下的智能程序到以 Internet 和 WWW 为平台的分布式智能系统；五是从智能技术的单一应用到各种各样的智能产品和智能工程（如智能交通、智能建筑）。

\q 参考答案要点：
（1）近期可能发展方向（3 分）：教材 9.7.3（印刷页 163）指出，未来人工智能可能会向 4 个方面发展：模糊处理、并行化、神经网络和机器情感。目前推理功能已获突破，学习及联想功能正在研究之中，下一步是模仿人类右脑的模糊处理功能和整个大脑的并行化处理功能；人工神经网络是未来人工智能应用的新领域；情感是智能的一部分，赋予计算机情感能力可能是下一个突破。
（2）发展趋势（3 分）：教材 9.7.3 还提出 4 个趋势：从独立进行的过程仿真走向与相关技术进行组合的功能仿真；从机器替代走向机器参与；从机器思维走向机器辅助人脑思维；从机器学习走向机器帮助人学习。
（3）可补充说明：教材 9.7.2（印刷页 162—163）指出，人工智能的研究与应用的结合越来越紧密，受应用的驱动越来越明显，商业上的成功成为实验室研究工作的催化剂；同时智能机器在意识、推理、创造力、直觉等方面不可能超越人类，应正确看待和使用智能机器人。

\answerpart{六、综合题（论述，1 题，共 16 分）}

\q 参考答案要点：
（1）八大类主要内容在农业中的对应（3 分）：搜索与求解——农业路径规划、农机调度、灌溉方案寻优；学习与发现——从历史产量、气象、土壤数据中发现规律；知识与推理——农业专家系统依据知识库进行病虫害诊断与施肥建议；发明与创造——育种方案、种植模式与农业软件创新；感知与交流——农业机器视觉识别、农业语音交互；记忆与联想——农业知识库与关联检索；系统与建造——智慧农业系统、农业机器人平台的设计实现；应用与工程——各类农业信息化工程项目的落地。（教材 9.4，印刷页 152—154）
（2）15 类应用在农业中的体现（4 分，答出 5 类即可，每类 0.8 分）：难题求解——病虫害诊断、产量预测等复杂问题；自动规划、调度与配置——农机作业路径规划与机群调度；机器定理证明——农业模型与规则的逻辑验证（教材方法的迁移应用）；机器翻译——农业技术资料的多语种翻译；智能控制——温室环境控制、精准灌溉控制；智能管理——农场生产经营管理系统；智能决策——农业生产与市场决策支持系统；智能通信——农业物联网通信网络的自适应管理；智能仿真——作物生长与农业系统仿真；智能人机接口——农业语音助手、可视化交互界面；模式识别——病虫害图像识别、果实分级、遥感图像解译；数据挖掘与知识发现——从农业大数据中发现增产规律；机器博弈——农业装备协同作业的博弈优化（较少见，可略）；智能机器人——采摘机器人、除草机器人、无人驾驶拖拉机。（教材 9.6，印刷页 155—161）
（3）智慧农业典型场景及支撑技术（4 分）：病虫害识别——机器视觉、深度学习、图像数据标注；产量预测——遥感、气象与土壤数据、机器学习回归与时间序列模型；智能农机——GPS/北斗导航、自动控制、机器视觉；精准灌溉——土壤墒情传感器、物联网、智能控制与优化算法。（大纲第五部分；教材第 9 章未展开农业应用，可结合教材第 10 章 10.2.3（7）农业自动化，印刷页 179—180，其中提到利用监测手段获取田间肥力、墒情、苗情、杂草、病虫害信息，自动进行精准耕作、播种、施肥、灌溉、除草、喷药、收割，即精准农业的核心内容）
（4）伦理与挑战（3 分，答出 3 点即可）：数据隐私与安全——农业经营数据、农户个人信息在采集与共享中面临泄露风险；算法偏见与公平——训练数据偏差可能导致识别错误或服务不均；就业替代——智能农机与农业机器人可能替代部分体力劳动，需重视农民转岗培训；安全可控与责任归属——自动驾驶农机、农业机器人出现故障或事故时的责任界定；数字鸿沟——小农户与规模化农场在技术获取和应用能力上的差距。教材 9.7.2（印刷页 162—163）强调机器人在意识、推理、创造力、直觉等方面不可能超越人类，应正确看待和使用智能机器人，不能盲目恐惧也不能放任不管。
（5）发展趋势与意义（2 分）：人工智能正向模糊处理、并行化、神经网络、机器情感方向发展，呈现从机器替代走向机器参与、从机器思维走向机器辅助人脑思维、从机器学习走向机器帮助人学习等趋势（教材 9.7.3）。对农业信息化而言，人工智能推动农业生产从经验驱动转向数据驱动，提高资源利用效率与生产精准度，是智慧农业和农业现代化的重要支撑；同时应坚持人机协同，重视伦理规范与人才培养。（结合 2021 年“人工智能与智慧农业”、2023 年“人工智能的应用”、2024 年“人工智能技术在农业上的应用”等真题）

\end{document}
